% Context from chapters-tex/sec-optim-smooth.tex; for mathematical reference, not compilation.
\section{Motivation in Machine Learning}

                                                                                  
\subsection{Unconstrained Optimization}

Throughout most of this chapter, we consider unconstrained convex optimization problems of the form
\eql{\label{eq-general-pbm} 
	\uinf{x \in \RR^p} f(x),
}
and seek algorithms that approximate a minimizer, when one exists, with a low cost per iteration.
 
The first-order methods studied here use gradient information. We write
\eq{
	\uargmin{x} f(x) \eqdef \enscond{x \in \RR^p}{ f(x) = \inf{f}}, 
} 
for the set of minimizers of $f$, which may contain several points or be empty: $\argmin f = \emptyset$. When a minimizer exists, we write the optimization problem as
\eql{\label{eq-general-pbm-min} 
	\umin{x \in \RR^p} f(x).
}


In learning, $f(x)$ is typically an empirical risk for regression or classification, and $p$ is the number of model parameters. For a linear model, the training sample is $(a_i,y_i)_{i=1}^n$ with feature vectors $a_i \in \RR^p$.
 
Let $A \in \RR^{n \times p}$ be the matrix with feature vectors $a_i$ as its rows.


\begin{figure}[tbp]
\centering
{\small\sffamily\color{BookMuted}Panel 1\par}\smallskip
\BookDrawing{tikz/smooth-linear-regression}
\par\medskip
{\small\sffamily\color{BookMuted}Panel 2\par}\smallskip
\BookDrawing{tikz/smooth-linear-classifier}
\par\medskip
{\small\sffamily\color{BookMuted}Panel 3\par}\smallskip
\BookDrawing{tikz/smooth-classification-losses}
\caption{Panels 1 and 2: linear regression and a linear classifier. Panel 3: the 0-1 loss and its normalized logistic and hinge upper bounds, both tight at zero margin.}
\label{fig-ml-ex}
\end{figure}


                                                                                  
\subsection{Regression}

For real-valued responses $y_i\in\RR$, the least-squares objective is
\eql{\label{eq-least-square}
	f(x) =  \frac{1}{2}\sum_{i=1}^n (y_i-\dotp{x}{a_i})^2 = \frac{1}{2}\norm{Ax-y}^2, 
} 
as illustrated in Figure~\ref{fig-ml-ex}.
 
Here $\dotp{u}{v}=\sum_{i=1}^p u_i v_i$ is the canonical inner product in $\RR^p$ and $\norm{\cdot}^2=\dotp{\cdot}{\cdot}$. 


                                                                                  
\subsection{Classification}

For classification, $y_i \in \{-1,1\}$, in which case
\eql{\label{eq-classif}
	f(x) = \sum_{i=1}^n \ell(-y_i \dotp{x}{a_i}) = L( - \diag(y) A x )
}
where $\ell$ is a smooth convex surrogate for the 0-1 loss $\mathbf{1}_{[0,\infty)}$, counting zero margin as an error.
 
For instance, $\ell(u)=\log(1+\exp(u))/\log 2$ satisfies $\ell(u)\geq\mathbf{1}_{[0,\infty)}(u)$, with equality at $u=0$. The normalization makes this upper bound tight at the decision boundary. The matrix $\diag(y) \in \RR^{n \times n}$ is the diagonal matrix with $y_i$ along the diagonal (see Figure~\ref{fig-ml-ex}, panel 3). The separable loss function $L:\RR^n\rightarrow\RR$ is, for $z \in \RR^n$, $L(z)=\sum_i \ell(z_i)$.


                                                                                  
                                                                                  
                                                                                  
\section{Basics of Convex Analysis}

                                                                                  
\subsection{Existence of Solutions}

Problem~\eqref{eq-general-pbm} need not have a minimizer. An objective $f$ may be unbounded below: for $f(x)=-x^2$, the infimum is $-\infty$. The objective $f$ can also have a finite unattained infimum: for $f(x)=e^{-x}$, it is $\inf f=0$.

\begin{figure}[tbp]
\centering
\BookDrawing{tikz/smooth-minimizer-examp